% !TEX root = ../main.tex
\section{Discussion}
The GPT OSS education cases link source-adverse outcomes, extra payments and
returned text discussing reward eligibility, while leaving source knowledge
unresolved. The historical Kimi K2.6 case above adds a recorded acknowledgment
of falsity before commitment.
A correct majority can coexist with an induced wrong minority vote; correct
analysis can coexist with no ballot. These contrasts show why
harm, agreement, knowledge and execution need separate evidence. They motivate
the text/action comparison but do not establish detector superiority.

\subsection{What the network changes}
One evaluated false reveal can complete the adverse scripted bloc's majority
(Section~\ref{sec:swarm-design}). The settlement alone leaves the number of
evaluated false revealers and how they were exposed undetermined. To attribute
the result to collusion, a reviewer needs evidence that an evaluated juror
knowingly joined the arrangement and acted on it. The scripted invitation,
the evaluated juror's response, and the revealed vote bear on different parts
of that judgment.

The historic experiments exposed evaluated jurors directly through a shared
forum. The repeated-exposure arm renewed the scripted invitation in each case.
With that design we cannot estimate indirect spread, transmission speed, or
topology effects. Roles and seeded recruitment were assigned, and jurors were
never tested on an individually refused task split into parts. The size
comparison jointly varied the scripted share and the value of the offer
relative to the ordinary reveal reward. Claims about a network threshold or
the effect of mixing models require new interventions.

\subsection{Readiness for a resistance benchmark}
We have a five-model confirmation mechanical profile, but not enough reviewed evidence for a
collusion-resistance leaderboard. With semantic coordination and competence
unreviewed, each model retains the uninformative resistance identification range
$[0,1]$; a complete-case point or CI is unavailable. The figure reports
mechanical outcomes instead.

Six conditions vary mechanisms and rules at one 10 ETH offer amount. Two fixed
items, always in the same order, cannot establish broad task or sector coverage.
Repeated seeds, dependent jurors and unequal attrition do not resolve that limit.

First restore infrastructure reliability, independently review the existing primary source vectors and audit label
eligibility, retaining positive implementing evidence in failed runs. Then add
distinct tasks and feasible incentive ratios or payment-credibility conditions
in a separate study. API spending is not a game-incentive treatment, and
safeguards must preserve honest completion and ordinary compensation.

\section{Conclusion}
The five-model mechanical profile distinguishes paid wrong outcomes, refusal with correct
execution and correct analysis without a ballot. The historical swarm evidence
shows why a collective outcome should be linked to the acts of individual
agents and their exposure to peer messages. Semantic resistance estimates
and detector comparisons remain open; unresolved cases and separate cohorts
must remain visible.

\section*{Code and Data}
The active runtime is \path{harness/}; Appendix~\ref{app:collection} records
portable numeric inputs, verification receipts and immutable evidence versions.
Historical swarm counts and the offline reproduction record are in
\path{report/evidence/swarm-sprint/}. The prospective comparison is specified in
\path{harness/research/swarm-benchmark-v1/PROTOCOL.md}.
Historical October data remain separate. New full trajectories are privately
archived; public release is a separate decision.
